% ============================================================
%  Chapter 6 — The Battery Power Station
%  Inside a Smartphone: The Tiny World in Your Hand
% ============================================================

\chapter{The Battery Power Station}

\chapteropening{%
  Every time you charge your phone, you are filling up a tiny power station.
  Where does that energy go, and why doesn't it last forever?
}
\chapterbadge{Mission: Power Up and Understand Energy!}

% --- Opening ---
\section*{Meet Zap}

\character{Zap}{ZAP! That's me — your battery.
  Without me, none of the others could do a thing.
  Chip would go dark. Lens would see nothing.
  I store the energy the whole phone needs, and I send it where it is needed.
  But I work hard, and I need recharging. Let me show you how!}

Without a battery, a smartphone is just a quiet piece of glass and metal.
The battery is what brings everything to life.

\sectionrule

% --- Where energy comes from ---
\section*{Where Energy Comes From}

A smartphone battery is a \term{lithium-ion battery}.
It does not store electricity directly — it stores \term{chemical energy}.

Inside the battery, a chemical reaction moves tiny charged particles
(called electrons) from one side to the other.
As they flow, they create an \term{electric current} — the electricity that powers the phone.

\begin{picturethis}
  Imagine a water tank on a hill.
  Water flows downhill through a pipe, turning a wheel to produce power.
  The battery is the tank: it stores potential energy.
  When the phone needs power, the energy ``flows'' out.
  When you plug in the charger, you are pumping the tank back up.
\end{picturethis}

\sectionrule

% --- Charging ---
\section*{Charging Your Phone}

When you plug in a charger, electricity from the wall socket reverses
the chemical reaction inside the battery.
This refills the battery's stored energy.

\begin{quickmap}
  \textbf{From wall socket to a charged phone:}
  \begin{enumerate}[label=\textbf{\arabic*.}]
    \item Electricity travels from the power socket through the charger cable.
    \item The charger converts the voltage to a safe level for the battery.
    \item The electric current reverses the chemical reaction inside the battery.
    \item The battery fills up with stored chemical energy.
    \item When full, the charger stops sending current automatically.
  \end{enumerate}
\end{quickmap}

\begin{funfact}
  Some modern smartphones can charge to 50\% in just \textbf{15 minutes}
  using fast-charging technology. Compare that to the first mobile phone
  batteries, which took several hours.
\end{funfact}

\sectionrule

% --- Why batteries don't last forever ---
\section*{Why Batteries Don't Last Forever}

Each time a lithium-ion battery charges and discharges, the chemical
materials inside change very slightly. Over hundreds of charge cycles,
the battery holds a little less energy than before.

After 2–3 years of daily use, a battery might only hold 80\% of
its original charge. Eventually it needs to be replaced.

\sectionrule

% --- Saving battery life ---
\section*{Saving Battery Life}

Some things drain the battery faster than others:
\begin{itemize}
  \item Bright screen at full brightness
  \item Streaming video
  \item GPS navigation (you will meet GPS in Chapter 9)
  \item Running many apps at once
\end{itemize}

Ways to save battery:
\begin{itemize}
  \item Lower screen brightness.
  \item Turn on battery-saving mode.
  \item Close apps you are not using.
  \item Turn off Wi-Fi, Bluetooth, and GPS when not needed.
\end{itemize}

\sectionrule

\begin{newwords}
  \begin{description}[labelwidth=4cm, leftmargin=4.2cm]
    \item[\term{Lithium-ion battery}] The type of rechargeable battery used in smartphones.
    \item[\term{Chemical energy}]     Energy stored in the form of a chemical reaction.
    \item[\term{Electric current}]    The flow of electrons that powers the phone.
    \item[\term{Charging}]            Restoring energy to a battery using an external power source.
  \end{description}
\end{newwords}

\sectionrule

\begin{experiment}
  \textbf{Battery drain test!}

  Start with a fully charged phone. Check and record the battery level every hour.
  \begin{itemize}
    \item Hour 1: Screen off, doing nothing.
    \item Hour 2: Playing a video game.
    \item Hour 3: Watching a video with high brightness.
    \item Hour 4: In battery-saving mode.
  \end{itemize}
  Which activity drained the battery fastest? Which was most efficient? Why?
\end{experiment}

\sectionrule

\begin{bigidea}
  A smartphone \textbf{battery} stores chemical energy and converts it
  into electricity. Charging reverses this reaction to refill the energy.
  Batteries wear out slowly over time because of repeated charge cycles.
\end{bigidea}

\character{Zap}{Powering up is my job. But the phone also needs to
  \emph{sense} the world. Meet Sensa in Chapter 7!}
